\newpage
\section{Collasso dei vertici}
\label{sec:collasso_dei_vertici}

Al termine della fase di parsing si dispone di una collezione di primitive (segmenti) indipendenti tra loro. Nei modelli CAD è abbastanza comune che in uno stesso punto siano presenti più vertici sovrapposti, appartenenti a segmenti diversi. È quindi necessario effettuare il collasso dei vertici, comunemente noto come \emph{vertex snapping}. Questa operazione serve sia a ottimizzare la computazione, riducendo il numero di vertici da elaborare, sia a soddisfare il requisito, necessario per le fasi successive, di rappresentare la geometria come un grafo con vertici univoci.

È importante notare che in questa fase vengono considerati solamente i segmenti relativi ai muri: porte e finestre saranno gestite in maniera diversa più avanti (Sezione \ref{sec:ricostruzione_varchi}).

L'obiettivo è quindi prendere l'insieme dei segmenti e produrre una nuova rappresentazione in cui i vertici estremi che distano meno di $\epsilon$ l'uno dall'altro vengono collassati in un unico vertice rappresentativo. Per risolvere questo problema è necessaria una struttura dati spaziale efficiente; si è scelto di utilizzare un \emph{hash spaziale}, discusso in dettaglio nella Sezione \ref{sec:strutture_dati_spaziali}.

In questa fase è fondamentale tarare in modo adeguato il parametro $\epsilon$: un valore troppo grande rischia di collassare erroneamente vertici distinti (ad esempio due angoli ravvicinati ma geometricamente reali), mentre un valore troppo piccolo lascia vertici duplicati che possono compromettere la costruzione del grafo nella fase successiva.

La classe che implementa questa struttura è mostrata nella Figura~\ref{fig:spatialhash-uml}:
\begin{figure}[H]
	\centering
	\begin{tikzpicture}
		\begin{class}[text width=9cm]{SpatialHash}{0,0}
			\attribute{- epsilon : Real}
			\attribute{- epsilonSq : Real}
			\attribute{- grid : Map<CellCoord, VertexId [*]>}
			\attribute{- vertices : Point [*]}
			\operation{+ SpatialHash(epsilon : Real = 10\^{-4})}
			\operation{+ snap(p : Point) : VertexId}
			\operation{+ findNearest(p : Point) : VertexId}
			\operation{- getCell(p : Point) : CellCoord}
		\end{class}
		
		\begin{class}[text width=3.5cm]{CellCoord}{-7,-8}
			\attribute{+ x : Integer}
			\attribute{+ y : Integer}
		\end{class}
		
		\begin{class}[text width=3cm]{VertexId}{7,-8}
			\attribute{+ value : Integer}
		\end{class}
	\end{tikzpicture}
	\caption{Diagramma UML della classe \texttt{SpatialHash}, che implementa la struttura dati spaziale usata per il collasso dei vertici.}
	\label{fig:spatialhash-uml}
\end{figure}

Il parametro $\epsilon$ definisce il raggio di ricerca $r$ e, di conseguenza, la dimensione della cella della griglia (Sezione \ref{subsec:hash_spaziale_teoria}). Il metodo principale è \texttt{snap()}: dato un punto $p$, cerca nella griglia $3 \times 3$ circostante un vertice già esistente che disti meno di $\epsilon$; se lo trova, restituisce l'id (\texttt{VertexId}) del vertice esistente, sul quale il nuovo punto viene collassato senza essere inserito nella griglia; in caso contrario, il punto viene inserito come nuovo vertice e ne viene restituito l'id appena assegnato.

Il metodo \texttt{find\_nearest()} restituisce invece l'id del vertice più vicino a $p$ all'interno dell'intero hash, indipendentemente dalla distanza; questo metodo verrà utilizzato più avanti nella ricostruzione dei varchi (Sezione \ref{sec:ricostruzione_varchi}).

Utilizzando $\epsilon = 10^{-2}$, il numero di vertici sui segmenti dei muri del caso di studio si riduce da 177 a 88, ovvero una riduzione di circa il 50\%. Riducendo ulteriormente il valore di $\epsilon$, la riduzione ottenuta non varia in modo significativo, a conferma che il valore scelto è già sufficiente a collassare correttamente tutti i vertici duplicati presenti nel modello.

Un esempio di utilizzo della struttura (Algoritmo~\ref{alg:build_edges}), nella costruzione dell'insieme di archi (\texttt{Edge}) a partire dai segmenti dei muri, è il seguente:
\begin{algorithm}[htbp]
	\caption{Costruzione degli archi a partire dai segmenti, con collasso dei vertici}
	\label{alg:build_edges}
	\begin{algorithmic}[1]
		\Require $segmentList$: lista dei segmenti estratti dal parsing
		\Require $snapEps$: tolleranza $\epsilon$ per il collasso dei vertici
		\Ensure $edges$: lista degli archi con vertici univoci
		
		\State $hash \gets$ \Call{SpatialHash}{$snapEps$}
		\State $edges \gets \emptyset$
		
		\For{\textbf{each} $seg \in segmentList$}
		\State $v_1 \gets hash.\Call{Snap}{seg.start}$
		\State $v_2 \gets hash.\Call{Snap}{seg.end}$
		\If{$v_1 \neq v_2$}
		\State $edges.\Call{Append}{(v_1, v_2, seg.layer)}$
		\EndIf
		\EndFor
		
		\State \Return $edges$
	\end{algorithmic}
\end{algorithm}

Il controllo \texttt{v1 != v2} evita di introdurre archi degeneri, ovvero archi di lunghezza nulla generati da segmenti i cui due estremi vengono collassati sullo stesso vertice.

La struttura di un arco (\texttt{Edge}) è molto simile alla struttura \texttt{Segment} introdotta nella Sezione \ref{sec:parsing}, con la differenza che, in questo caso, vengono memorizzati solo i riferimenti (id) ai vertici anziché le loro coordinate (Figura~\ref{fig:edge-uml}):
\begin{figure}[H]
	\centering
	\begin{tikzpicture}
		\begin{class}[text width=5cm]{Edge}{0,0}
			\attribute{+ v1 : VertexId}
			\attribute{+ v2 : VertexId}
			\attribute{+ layer : SegmentLayer}
		\end{class}
	\end{tikzpicture}
	\caption{Diagramma UML della struttura \texttt{Edge}, che rappresenta un lato del grafo con la relativa semantica.}
	\label{fig:edge-uml}
\end{figure}

Al termine di questa fase si dispone quindi di un vettore di \texttt{Edge} e di un vettore di vertici univoci.